\documentclass[../DoAn.tex]{subfiles}
\begin{document}

\begin{center}
    \Large{\textbf{TÓM TẮT NỘI DUNG ĐỒ ÁN}}\\
\end{center}
\vspace{1cm}
% Sinh viên viết tóm tắt ĐATN của mình trong mục này, với 200 đến 350 từ. Theo trình tự, các nội dung tóm tắt cần có: (i) Giới thiệu vấn đề (tại sao có vấn đề đó, hiện tại được giải quyết chưa, có những hướng tiếp cận nào, các hướng này giải quyết như thế nào, hạn chế là gì), (ii) Hướng tiếp cận sinh viên lựa chọn là gì, vì sao chọn hướng đó, (iii) Tổng quan giải pháp của sinh viên theo hướng tiếp cận đã chọn, và (iv) Đóng góp chính của ĐATN là gì, kết quả đạt được sau cùng là gì. Sinh viên cần viết thành đoạn văn, không được viết ý hoặc gạch đầu dòng.
Trong thời đại số hóa hiện nay, thương mại điện tử đang phát triển mạnh mẽ và trở thành xu hướng phổ biến trong lĩnh vực kinh doanh. Tuy nhiên, phần lớn các giải pháp hiện tại tập trung vào các nền tảng web, trong khi nhu cầu mua sắm qua thiết bị di động ngày càng tăng cao. Bên cạnh đó, các hệ thống thương mại điện tử hiện có thường có cấu trúc phức tạp, khó mở rộng và bảo trì, đặc biệt là khi triển khai cả ứng dụng di động và hệ thống quản trị web song song. Một số hướng tiếp cận như xây dựng native app cho từng nền tảng hoặc sử dụng hệ thống quản lý riêng biệt cho web và mobile đều tồn tại những hạn chế như tốn kém chi phí phát triển, khó đồng bộ dữ liệu và khó kiểm soát hệ thống toàn diện.

Để giải quyết vấn đề trên, em lựa chọn hướng tiếp cận xây dựng hệ thống thương mại điện tử đa nền tảng sử dụng Flutter cho ứng dụng di động, React cho giao diện quản trị Web, Spring Boot cho Backend, và Firebase làm cơ sở dữ liệu chung. Hướng tiếp cận này giúp đảm bảo hiệu năng, khả năng mở rộng, dễ dàng bảo trì cũng như đồng bộ dữ liệu giữa các nền tảng. Flutter được lựa chọn nhờ khả năng phát triển ứng dụng đa nền tảng từ một codebase duy nhất, còn MVC giúp tách biệt logic, giao diện và xử lý dữ liệu rõ ràng.Tất cả dữ liệu được lưu trữ và đồng bộ qua Firebase, đảm bảo tính nhất quán và thời gian thực, cấu trúc NoSQL linh hoạt và tối ưu hóa hiệu năng cho mobile

Đóng góp chính của đồ án là xây dựng một hệ sinh thái thương mại điện tử đầy đủ, có khả năng hoạt động hiệu quả trên nhiều nền tảng với chi phí phát triển tối ưu. Kết quả cuối cùng là một hệ thống hoàn chỉnh, có tính thực tiễn cao, dễ dàng triển khai và mở rộng trong môi trường kinh doanh thực tế.
\end{document}